%! TEX root = gambar potongan.tex
% the main TeX file which is intended to compile, :VimtexReload after adjustment
   % this file should be included in the main file
% :h vimtex-tex-root
\documentclass[../menggambarteknik.tex]{subfiles}
%ensure the class of docs

\begin{document}

\section{Gambar Potongan}

\subsection*{Tujuan}
Setelah mempelajari bahan dalam bab ini, seharusnya Anda dapat:
\begin{enumerate}
	\item Menunjukkan macam-macam penyajian potongan.
	\item Mendemonstarasikan cara menggambar potongan.
	\item Mengklarifikasi gambar benda yang tidak boleh dipotong.
	\item Mendesain gambar beberapa benda yang dipotong.
\end{enumerate}

\subsection{Macam-Macam Gambar Potongan}
Penggunaan garis strip-strip (gores) untuk melukiskan bagian benda yang tidak terlihat dalam jumlah yang sedikit memang bisa membantu para pembaca gambar, tetapi bila bagian yang tidak terlihat banyak akan membingungkan. Untuk menghindari kebingungan dan memperjelas bagian dalam suatu benda yang akan digambar dipergunakan gambar potongan (sectional views).

Untuk memperlihatkan bagian dalam suatu benda dengan menggunakan gambar potongan dapat dilakukan dengan potongan seluruhnya, potongan separoh dan potongan sebagian disesuaikan dengan kadar kebutuhan dari bagian dalam yang akan diperlihatkan (lihat Gambar 31). Memang penggunaan gambar potongan seluruhnya akan lebih memperlihatkan bagian dalam, tetapi dalam hal-hal tertentu justru akan mubazir terutama dalam penggunaan waktu menggambar, seperti benda kerja yang simetris, maka gambar potongannya cukup separoh atau sebagian saja tidak perlu seluruhnya.

\begin{figure}
	\begin{center}
		\includegraphics[width=.95\textwidth]{macam potongan.jpg}
	\end{center}
	\caption{Macam-macam potongan}
	\label{fig:macam potongan}
\end{figure}


\subsection{Cara Menggambar Potongan}
Bagian dalam yang mendapat potongan perlu dibedakan dengan bagian luar yang tidak dipotong. Untuk itu seluruh bagian yang dipotong diarsir dengan sudut 45o terhadap garis sumbu atau garis gambar (lihat Gambar 32). Jarak garis arsir yang dibuat disesuaikan dengan besarnya gambar dan jaraknya sama antara satu sama lainnya. Gambar susunan benda kerja yang menjadi satu, potongannya ditunjukkan dengan arsiran yang berbeda arah (lihat Gambar 33), sedangkan potongan dari satu benda harus diarsir dengan arah yang sama. Untuk benda yang tipis gambar potongannya ditunjukkan tidak dengan arsir, tetapi cukup ditebalkan dengan warna hitam (lihat Gambar 34).

\begin{figure}
	\begin{center}
		\includegraphics[width=.95\textwidth]{arsir potongan.jpg}
	\end{center}
	\caption{Arsir untuk penunjukkan potongan}
	\label{fig:arsir potongan}
\end{figure}

\begin{figure}
	\begin{center}
		\includegraphics[width=.95\textwidth]{arsir beda susunan.jpg}
	\end{center}
	\caption{}
	\label{fig:Arsir berbeda untuk gambar susunan}
\end{figure}


Untuk penampang yang tipis, seperti benda yang terbuat dari plat, baja profil, dan paking dapat digambar dengan garis tebal (dihitamkan), sedangkan daerah yang dihitamkan dari beberapa penampang yang berbeda dipisahkan (diberi jarak) sedikit (lihat Gambar 34).

\begin{figure}
	\begin{center}
		
\includegraphics[width=.95\textwidth]{arsir hitam.png}
	\end{center}
	\caption{Arsir hitam untuk benda tipis}
	\label{fig:arsirhitam}
\end{figure}





Gambar potongan memang dapat membantu untuk menjelaskan bentuk bagian dalam yang tidak terlihat dan ukuran-ukurannya, namun demikian ternyata tidak semua benda dapat dipotong. Bagian-bagian benda seperti baut, paku keling, pasak, poros, dan sirip tidak boleh dipotong (lihat Gambar 35).

\begin{figure}
	\begin{center}
		\includegraphics[width=.95\textwidth]{benda tdk dipotong.jpg}
	\end{center}
	\caption{Benda-benda yang tidak boleh dipotong}
	\label{fig:bendatidakpotong}
\end{figure}


Pemotongan pada suatu pandangan dilakukan dengan menggunakan garis potong, yaitu garis strip titik dengan ujung tebal dan diberi anak panah yang diberi huruf sama. Pada penunjukkan bagian yang dipotong ditulis huruf yang sama dengan pemotongannya (lihat Gambar 36).

\begin{figure}
	\begin{center}
		\includegraphics[width=.95\textwidth]{penunjukkan pemotongan.png}
	\end{center}
	\caption{Penunjukkan Pemotongan}
	\label{fig:pemotongan}
\end{figure}




\bibliographystyle{apalike}
\bibliography{../filename.bib} % required for citecompletion, not work in mainfile
\end{document}

